% Modernised reading — fonds Alexandre Grothendieck, shelfmark 111, pages 1-16
% (the whole folder).
%
% This is an INTERPRETATION, not a transcription. It re-reads the pages in
% today's notation and vocabulary, states the hypotheses the typescript leaves
% implicit, and drops the critical apparatus so that the mathematics reads
% straight through; everything the apparatus carried moves into footnotes. It is
% held to being mathematically correct as it stands: where the pages are loose
% or mistaken, the true statement is what appears here and a footnote says what
% the page has. In any dispute of substance the transcription governs: whoever
% cites Grothendieck must cite batch-01.fr.tex (pp. 1-16, the whole folder).
%
% Pass: Opus 5.5 (claude-opus-5-5), 2026-10-03 — first pass, unchecked against the pages by a human.
% The folder was read whole (one batch); this is one document for the
% shelfmark. It was made on Opus 5.5 and has not been measured against the
% reference reading of folder 115 (made on Opus 5).
%
% Witness. 111 is a photocopy of the sixteen-page typescript « Tapis de
% Quillen » (notes du 10.9.68): typescript page N = page N of this folder. Two
% more witnesses are in folder 162-5: copy B (162-5 pp. 32-41, 43-48, the
% annotated top copy) and copy A (162-5 pp. 2-17, a photocopy of B after the
% ink, with marks of its own). 111 reproduces B's ink and A's own marks (batch
% header, « Relation to folder 162-5 »). This reading follows 111's
% transcription throughout; copy B is invoked only where 111's transcription
% itself cites it (pp. 7, 10, 13, 16) or where a footnote says so explicitly.
% It is written to agree in substance with 162-5.modern.tex, whose departures
% it carries over with 111's page numbers.
%
% Notation fixed for the whole document (as in 162-5.modern.tex). N(C) for the
% nerve (the page's S(C)); Ho(Cat), Ho(sSet) for (Cat)', (Ssimpl)'; C^ for the
% presheaf topos (the page's C~, once C^); « catégorie double » for the page's
% « bi-catégorie »; V for the category of manifolds of part III (the page's C);
% E for the target category of the universal problem (the page's B, kept bold
% \mathbf B for Quillen's solution); B_d(X) graded by relative dimension;
% H^k(X) = H^k(X; F_2), matched with B_{-k}; W for the page's T-bar' (p. 13);
% Hom_{B^H} for the page's H(X,Y).
%
% Substantive departures, with pages.
%  p. 1   coefficient systems = functors C° -> Ens inverting EVERY arrow (the
%         page: « transformant isomorphismes en isomorphismes »); lim^(i)_C F
%         read as lim^i over C° of u*F'.
%  p. 3   the typed coproduct of cochains corrected to a product (the hand's
%         « ∏ ! »).
%  p. 3   the converse of the D^b_lc statement holds over Z, not over any
%         non-zero ring (counterexample K(Z/p,2), ours).
%  p. 4   C, C' read X, X'; « se borner aux systèmes locaux commutatifs »
%         answered negatively (acyclic groups, ours).
%  p. 5   unit functor supplied; Segal condition supplied for « objets qui … »;
%         « un triple » for a five-term datum.
%  p. 7   Hom_grab identified with the Picard category of tau^{<=0} RHom
%         (truncation supplied); the struck commutative case explained.
%  p. 8   « pl. fidèle (?) »: the loop functor is not even faithful (ours);
%         the « tuage » analogy for Aut(1) footnoted as rather a looping.
%  p. 11  compact (closed) manifolds supplied (the page's own « (compactes ?) »).
%  p. 12  F^• : V° -> B (the page: C -> B); composition in B written out;
%         « au dessus de X » read « au-dessus de X x Y ».
%  p. 13  base change written alpha_{s•} alpha_s^• = Gamma^• i_{s•} (the page's
%         order alpha_s^* alpha_{s*} does not compose).
%  p. 13  H (x) Lambda = B holds as Lambda-modules for any homogeneous lift
%         (Conner-Floyd freeness, graded Nakayama).
%  p. 13-14 « quasi-abélienne » = pseudo-abelian (Karoubian additive), not
%         Schneiders' quasi-abelian.
%  p. 16  « Thom's theorem » (x in S(X) iff A^+ x = 0), read with the usual
%         action on cohomology, fails: RP^1 in RP^2. Reported under reserve.
%
% Announced and never established on the pages: the generalisation to
% « localement homotopiquement triviaux » topoi (pp. 3-4); n-categories as
% n-simplicial sets (p. 6); the Hom_grab identification (p. 7) and the
% 2-Grab-category of homomorphisms (p. 9); abelian groups in n-Cat <->
% complexes of length n (p. 9); K_n via universal Gr-n-categories (p. 9);
% the formalisation for general correspondences (p. 12); « le reste » of the
% universal property, « sans doute l'AQT » (p. 13); the universality of S
% (p. 15); the structure of G (p. 16).
%
% Dating inference (not a reading): A's own mark « cf exposé Karoubi à
% Bourbaki », reproduced on 111 p. 16, most likely points to Karoubi's
% Bourbaki talk of 1971-72; if so, the photocopy 111 was made after 1971.
% \dating keeps the inventory's « 1968 », the date of the typed text.
\documentclass[11pt,a4paper]{article}
\input{../preamble/grothendieck}

\folder{111}
\pages{1}{16}
\dating{1968}
\watermark{Édition de démonstration\\interprétation personnelle de l'œuvre}
\foldertitle{Tapis de Quillen [notes du 10/09/1968] : copies de tapuscrit annoté (1968) — lecture modernisée du dossier entier}

\begin{document}

\section*{Résumé}

\begin{resume}
Ces seize pages dactylographiées sont des notes prises en septembre 1968 sur ce
que Daniel Quillen, alors âgé de vingt-huit ans, exposait à Grothendieck. Le
dossier n'en contient qu'une photocopie, qui reproduit les annotations
manuscrites portées sur un autre exemplaire ; celui-ci et une seconde copie se
trouvent au dossier 162-5. « Tapis », dans la
langue de ces années, désigne un ensemble de résultats et la manière de s'en
servir : le texte parle aussi du « tapis de Dold » et du « tapis de Roos ».

Le tapuscrit a trois parties. La première part de ce qu'une catégorie est un
espace : ses objets sont des points, ses flèches des segments, ses triangles
commutatifs des faces. Selon Quillen, les catégories et les espaces
combinatoires qu'on appelle ensembles simpliciaux ont la même théorie de
l'homotopie, et toute catégorie a le même type d'homotopie que son opposée,
celle où l'on a retourné toutes les flèches. Le texte y voit l'occasion de
définir un type d'homotopie par ses « systèmes locaux », c'est-à-dire par la
façon dont des coefficients peuvent tourner le long des lacets de l'espace.

La deuxième partie traite des catégories munies d'une loi de groupe. Un groupe
ordinaire exige que $(ab)c$ soit \emph{égal} à $a(bc)$ ; dans une catégorie, on
peut se contenter d'un isomorphisme entre les deux, pourvu qu'il soit donné et
cohérent. Une telle « catégorie de Picard » se décrit, à peu de chose près, par
une flèche entre deux groupes abéliens, un complexe de longueur un. Il en va de
même en longueur deux, et le texte espère qu'il en ira ainsi en toute longueur.
Il espère aussi définir des invariants $K_n$ d'une catégorie munie d'un produit
comme les groupes d'homotopie de sa « complétion en groupe », comme on fabrique
les entiers relatifs à partir des entiers naturels. C'est ce que Quillen fera
dans les années suivantes, par d'autres chemins.

La troisième partie regarde le cobordisme d'un œil de géomètre algébriste. Deux
variétés sont cobordantes quand, ensemble, elles bordent une troisième, comme
deux cercles bordent un cylindre. En comptant les variétés au-dessus d'un espace
à cobordisme près, on obtient une catégorie de « correspondances » qui ressemble
à la catégorie des motifs, cette théorie universelle des cohomologies que
Grothendieck développait alors pour les variétés algébriques. Le groupe des
symétries des opérations de Steenrod y joue le rôle d'un « groupe de Galois
motivique ». Contrairement au cas familier, ce groupe est unipotent, et il est
défini sur le corps à deux éléments.

Ces pages montrent aussi une manière de travailler. Le texte rapporte
(« Quillen prouve ») et se pose des questions (« Demander à Quillen… »). Dans les
marges, la main répond à certaines d'entre elles, par « oui », par « Si ! » ou
par « non, … ». On y lit une conversation mathématique en train de se faire :
les choses comprises, celles qu'on croit, celles qu'on ira demander. La
photocopie porte en outre quelques notes plus tardives, dont un renvoi à un
exposé de Karoubi qui laisse penser qu'elle a été tirée quelques années après
1968. Le dossier est classé parmi les pièces rassemblées autour de « À la
poursuite des champs », le texte de 1983 qui commence lui-même par une lettre à
Quillen.

Pour aller plus loin, les noms à chercher sont : la structure de modèles de
Thomason sur les catégories ; la correspondance de Dold--Kan ; les 2-groupes et
les champs de Picard ; la $K$-théorie supérieure de Quillen ; et surtout la note
de Quillen de 1969 sur les lois de groupes formels du cobordisme, qui donnera à
la troisième partie sa structure véritable, et que prolonge aujourd'hui le
cobordisme algébrique de Levine et Morel.

\keywords{nerve of a category, category of simplices, homotopy theory of categories, locally constant sheaf, double category, Segal condition, Dold–Kan correspondence, Picard category, 2-group, higher algebraic K-theory, group completion, unoriented cobordism, correspondence category, Karoubi envelope, Chow motive, Steenrod algebra, motivic Galois group, Pontryagin–Thom construction}
\end{resume}

\pagerange{1}{16}
\section*{Le fil du dossier, et les conventions}

\emph{La pièce.} Le dossier est la photocopie d'un tapuscrit de seize pages,
numérotées 1 à 16 à la main en tête de chaque page. La page $N$ du tapuscrit
est la page $N$ du dossier. Le texte est titré « Tapis de Quillen. » ; la main
a daté la première page « notes du 10.9.68 », d'où le « [notes du 10/09/1968] »
de l'inventaire. La machine n'avait ni lettres grecques ni flèches ; plusieurs
ont été tracées ensuite à la main, avec des corrections de notation et des
notes de fond.

\emph{Les trois exemplaires.} Le dossier 162-5 contient deux autres témoins du
même original dactylographié, avec les mêmes coupures de lignes et les mêmes
biffures à la machine. L'un est l'original annoté à l'encre bleue (162-5, pages
32 à 48), que nous appelons l'\emph{exemplaire annoté}. L'autre est une
photocopie noire de celui-ci, tirée après l'encre et qui a reçu quelques notes
propres (162-5, pages 2 à 17). La couche manuscrite de 111 est celle de
l'exemplaire annoté, photocopiée, plus ces notes propres, reproduites des mêmes
traits aux mêmes places. On ne trouve sur 111 aucune marque qui ne soit sur
l'un des deux autres. 111 n'est pourtant pas simplement une copie de la
photocopie de 162-5 telle qu'elle est aujourd'hui : il garde au bord droit de
quelques pages (7 et 13) un texte que celle-ci a perdu. Les images ne disent
pas s'il a été tiré sur elle avant cette perte, ou sur un intermédiaire
commun.\footnote{Tout ceci est le constat de la transcription de 111, établie
puis revue sur l'exemplaire annoté, et non une lecture nouvelle. Le texte suivi
ici est partout celui de 111. L'exemplaire annoté n'est invoqué qu'aux endroits
où la transcription de 111 le cite elle-même pour lire un passage pâli par la
photocopie, ou quand une note le dit.}

\emph{Qui écrit.} Le tapuscrit parle à la première personne : « ma lettre à
Verdier », « “ma” construction chérie de la “catégorie de Picard” », « mon
interprétation chérie », « le rôle universel que j'ai envisagé plus haut ». Il
rapporte ce que « Quillen prouve » et note ce qu'il faudra « demander à
Quillen ». Rien sur les pages ne le signe. Une vérification de l'encre sur le
fac-similé, faite par sondage à la page 3, tient les ajouts manuscrits pour
siens ; la transcription ne va pas plus loin, et nous non plus.

\emph{Le fil.} Les trois parties du tapuscrit ne s'enchaînent pas par
déduction ; elles sont trois sujets d'une même conversation.
\begin{itemize}
  \item \emph{I, pages 1 à 5} : catégories et ensembles simpliciaux. Le nerf,
  la catégorie des simplexes, l'équivalence des théories homotopiques, $C$ et
  $C^{\circ}$, puis une généralisation proposée aux topos « localement
  homotopiquement triviaux » et l'idée de définir un type d'homotopie par la
  catégorie dérivée de ses systèmes locaux.
  \item \emph{II, pages 5 à 10} : catégories doubles, 2-catégories, groupes
  dans $(\mathrm{Cat})$ et dans $(2\text{-Cat})$, la correspondance de
  Dold--Kan, et le programme des $K_n$ d'une catégorie monoïdale.
  \item \emph{III, pages 11 à 16} : le cobordisme comme théorie de
  correspondances, la catégorie universelle $\mathbf{B}$, sa description par
  des $\Lambda$-modules, l'analogie avec les motifs, une seconde catégorie
  $\mathbf{S}$ et le groupe $G$ dual de l'algèbre de Steenrod, enfin la
  généralisation orientée.
\end{itemize}
Une dernière section revient sur la couche manuscrite de la photocopie, et sur
ce qu'elle laisse inférer de sa date.

\emph{Notations}, valables pour tout le document ; elles sont celles de la
lecture modernisée du dossier 162-5, et chaque écart avec la page est signalé
en note. $N(C)$ désigne le nerf (la page : $S(C)$) et $C^{\circ}$ la catégorie
opposée. $\mathrm{Ho}(\mathrm{Cat})$ et $\mathrm{Ho}(\mathrm{sSet})$ désignent
les catégories homotopiques (la page : $(\mathrm{Cat})'$, $(\mathrm{Ssimpl})'$),
et $\widehat{C}$ le topos des préfaisceaux sur $C$ (la page : $\widetilde{C}$).
« Catégorie double » rend la « bi-catégorie » de la page. Dans la troisième
partie, $\mathcal{V}$ est la catégorie des variétés (la page : $C$),
$\mathcal{E}$ la catégorie but du problème universel (la page : $B$), et
$H^{\bullet}(X) = H^{*}(X; \mathbb{F}_2)$.

\emph{Annoncé, non établi.} Le tapuscrit laisse ouverte la généralisation aux
topos localement homotopiquement triviaux (pages 3 et 4). Il pose sans les
établir la description des $n$-catégories par des ensembles $n$-simpliciaux
(page 6) et celle des homomorphismes entre catégories de Picard (pages 7 et 9).
Il ne fait qu'espérer la correspondance entre groupes abéliens de
$(n\text{-Cat})$ et complexes de longueur $n$, et la définition des $K_n$
(page 9). Il renvoie « le reste » de la propriété universelle de $\mathbf{B}$ à
une abréviation que nous ne savons pas lire (page 13), suppose sans la
démontrer l'universalité de la catégorie $\mathbf{S}$ (page 15), et ne dit pas
la structure du groupe $G$ (page 16).

\pagerange{1}{5}
\section*{I. Catégories et ensembles simpliciaux (pages 1 à 5)}

\pagerange{1}{2}
\subsection*{Le nerf, le topos, la catégorie des simplexes (pages 1 et 2)}

\emph{Le nerf.} À toute catégorie $C$ on associe son nerf $N(C)$, l'ensemble
simplicial des suites de flèches composables. Le foncteur $N : (\mathrm{Cat})
\to (\mathrm{sSet})$ est pleinement fidèle.\footnote{La page note $S(C)$ le nerf
et $(\mathrm{Ssimpl})$ la catégorie des ensembles semi-simpliciaux,
c'est-à-dire, dans la terminologie de l'époque, simpliciaux. Nous réservons la
lettre $S$ à la troisième partie.} Les systèmes locaux d'ensembles sur $N(C)$
correspondent aux foncteurs sur $C$ qui envoient toute flèche sur un
isomorphisme, c'est-à-dire qui se factorisent par le groupoïde enveloppant
$C[C^{-1}]$. Leur cohomologie s'interprète de deux façons : par les foncteurs
dérivés de $\varprojlim$, ou comme cohomologie du topos $\widehat{C}$ des
préfaisceaux sur $C$. Il s'agit de $H^{0}$ pour les ensembles, de $H^{1}$ pour
les groupes, de tous les $H^{i}$ pour les groupes abéliens. Pour un système
local abélien $F$, vu comme foncteur $C^{\circ} \to \mathrm{Ab}$, on a
$H^{i}(N(C), F) = H^{i}(\widehat{C}, F) = \varprojlim^{i}_{C^{\circ}} F$.

Un foncteur $u : C \to C'$ est un \emph{homotopisme} si $N(u)$ est une
équivalence faible ; c'est le terme de la page. Le critère cohomologique
d'Artin et Mazur le caractérise : $u$ est un homotopisme si et seulement si,
pour tout système de coefficients local $F'$ sur $C'$, les applications
\[
\varprojlim\nolimits^{i}_{C'^{\circ}} F' \longrightarrow
\varprojlim\nolimits^{i}_{C^{\circ}} u^{*}F'
\]
sont bijectives, pour les $i$ où cela a un sens.\footnote{La page écrit au
second membre $\varprojlim^{(i)}_{C} F$, au lieu de l'image inverse de $F'$, et
indexe $\varprojlim$ par $C$ plutôt que par $C^{\circ}$. La parenthèse ouverte
devant « du \emph{topos} » n'est pas refermée, et ce $\widehat{C}$ porte un
chapeau pâle là où le reste de la page trace des tildes à la main.
« Homotopisme » se dit aujourd'hui équivalence faible de catégories. Le
théorème A de Quillen (1973) en donnera un critère suffisant par les catégories
comma ; le livre d'Artin et Mazur, « Etale homotopy », paraîtra en 1969. Les
deux sont postérieurs à ces notes.}

\emph{Le topos.} À $C$ est associé le topos $\widehat{C}$, qui varie de façon
\emph{covariante} avec $C$ : un foncteur $u$ induit un morphisme de topos
$\widehat{u} : \widehat{C} \to \widehat{C'}$, dont l'image inverse est la
restriction le long de $u$. La page ajoute que le foncteur $C \mapsto
\widehat{C}$ n'a plus rien de pleinement fidèle, « semble-t-il ?? ».\footnote{C'est bien le cas : le topos des préfaisceaux ne voit $C$ qu'à
complétion de Karoubi près, et deux catégories de même enveloppe karoubienne
ont des topos de préfaisceaux équivalents. Les tildes de la page sur les $C$
de ce paragraphe et du suivant sont tracés à la main.} Un système de
coefficients ensembliste sur $C$ est un foncteur $C^{\circ} \to \mathrm{Ens}$
qui envoie toute flèche sur une bijection.\footnote{La page les définit comme
les foncteurs « transformant isomorphismes en isomorphismes », ce que fait tout
foncteur. Le sens est celui du premier alinéa : toute flèche est envoyée sur un
isomorphisme.} Ce sont exactement les objets localement constants de
$\widehat{C}$, qui se définissent intrinsèquement à partir du topos. Il en
résulte qu'être un homotopisme ne dépend que de $\widehat{u}$ : cela signifie
que $H^{i}(\widehat{C'}, F') \to H^{i}(\widehat{C}, \widehat{u}^{*}F')$ est
bijectif pour tout faisceau localement constant $F'$ sur $\widehat{C'}$.

\emph{La catégorie des simplexes.} Dans l'autre sens, on associe à un ensemble
simplicial $X$ la catégorie de ses simplexes $T(X) = \Delta/X$. Ses objets sont
les simplexes de $X$ ; c'est une catégorie fibrée sur $\Delta$, à fibres les
ensembles discrets $X_n$. Quillen prouve deux choses. D'abord, $N T(X)$ est
canoniquement isomorphe à $X$ dans $\mathrm{Ho}(\mathrm{sSet})$. Ensuite,
$T N(C)$ est canoniquement isomorphe à $C$ dans la localisée
$\mathrm{Ho}(\mathrm{Cat})$ de $(\mathrm{Cat})$ par les homotopismes. Ces
isomorphismes sont fonctoriels. Il en résulte formellement qu'un morphisme $f$
d'ensembles simpliciaux est une équivalence faible si et seulement si $T(f)$
est un homotopisme, et que $N$ et $T$ induisent des équivalences
quasi-inverses
\[
\mathrm{Ho}(\mathrm{Cat}) \ \simeq\ \mathrm{Ho}(\mathrm{sSet}).
\]
C'est aujourd'hui un résultat classique.\footnote{On le cite d'ordinaire
d'après Illusie (1972), qui l'attribue à Quillen ; Thomason (1980) en fera une
équivalence de catégories de modèles. L'isomorphisme $N(\Delta/X) \to X$ est
réalisé par l'application « dernier sommet ».}

\pagerange{2}{3}
\subsection*{Une catégorie et son opposée (pages 2 et 3)}

Quillen construit de plus un isomorphisme canonique et fonctoriel entre $C$ et
$C^{\circ}$ dans $\mathrm{Ho}(\mathrm{Cat})$, ou, ce qui revient au même,
entre $N(C)$ et $N(C^{\circ})$ dans $\mathrm{Ho}(\mathrm{sSet})$. C'est un
isomorphisme de la catégorie localisée, non un foncteur de $C$ dans
$C^{\circ}$.\footnote{La page ne dit pas comment Quillen le construit, et nous
n'en disons pas plus qu'elle.} Son effet sur les systèmes locaux est le
suivant : un foncteur contravariant $F$ sur $C$, qui rend toute flèche
inversible, devient le foncteur covariant qui a les mêmes valeurs sur les
objets et $F(u)^{-1}$ sur une flèche $u$. Sur les groupoïdes fondamentaux,
c'est l'isomorphisme évident $\pi_1(C) \to \pi_1(C^{\circ}) =
\pi_1(C)^{\circ}$, $g \mapsto g^{-1}$.

Il s'en déduit une interprétation faisceautique de la cohomologie d'un ensemble
simplicial $X$ à coefficients dans un système local covariant $F$, définie
classiquement par le complexe cosimplicial
\[
C^{n}(F) = \prod_{x \in X_n} F(x).
\]
On passe au système local contravariant défini par $F$, on le regarde comme un
faisceau sur $T(X)$, c'est-à-dire comme un objet de $(\mathrm{sSet})_{/X}$, et
on prend sa cohomologie.\footnote{La frappe porte $\coprod_{x \in X_n}$ ; la
main écrit en regard un signe de produit suivi d'un point d'exclamation. Des
cochaînes sont bien un produit, et nous écrivons $\prod$. « covariant » est une
lecture incertaine : la frappe porte « con », trois lettres biffées à la
machine, puis « variant », de sorte que « contravariant » n'est pas exclu. La
suite, qui passe au système « contravariant défini par $F$ », appelle
« covariant ».} La page ajoute que, pour un système de coefficients covariant
non local, on n'a toujours pas d'interprétation faisceautique de la cohomologie
classique ; ni, pour un système contravariant, de son homologie, ni
inversement de sa cohomologie faisceautique en termes classiques. Un grand
crochet manuscrit embrasse ce paragraphe, et la marge répond :
« Si ! ».\footnote{La réponse ne dit pas laquelle des trois lacunes elle
comble. Pour la première, il y a bien une réponse : un système de coefficients
covariant quelconque sur $X$ est un faisceau sur la catégorie opposée
$T(X)^{\circ}$, et sa cohomologie classique est celle des $\varprojlim^{i}$
sur $T(X)$. Cette remarque est la nôtre.}

\pagerange{3}{5}
\subsection*{Catégories dérivées de systèmes locaux, et topos localement homotopiquement triviaux (pages 3 à 5)}

Soit $u : C \to C'$ un homotopisme. Notons $D^{b}_{\mathrm{lc}}(C)$ la
sous-catégorie triangulée de la catégorie dérivée bornée des faisceaux abéliens
sur $\widehat{C}$ formée des complexes dont les faisceaux de cohomologie sont
localement constants. Quillen montre que $u^{*} : D^{b}_{\mathrm{lc}}(C') \to
D^{b}_{\mathrm{lc}}(C)$ est une équivalence, et réciproquement. La partie
directe vaut avec un anneau de base quelconque, et aussi avec un anneau de
coefficients constant tordu. La réciproque vaut à coefficients dans
$\mathbb{Z}$, mais pas dans un anneau quelconque non nul.\footnote{La page dit
« n'importe quel anneau de base (à condition de le supposer $\neq 0$ dans le
cas de la réciproque) ». Pour la réciproque, c'est faux. Soit $C$ une catégorie
dont le nerf est un $K(\mathbb{Z}/p, 2)$, et $u : C \to e$. À coefficients dans
$\mathbb{Q}$, ou dans $\mathbb{F}_{\ell}$ avec $\ell \neq p$, les systèmes
locaux sur $C$ sont constants, puisque $C$ est simplement connexe, et
$H^{i}(C, M) = 0$ pour $i > 0$ ; $u^{*}$ est donc une équivalence, et $u$
n'est pas un homotopisme. À coefficients dans $\mathbb{Z}$, l'équivalence des
systèmes locaux donne l'isomorphisme des groupes fondamentaux, et celle des
cohomologies à coefficients locaux donne l'homotopisme, par le théorème de
Whitehead. Le contre-exemple est le nôtre ; c'est aussi celui de la lecture
modernisée du dossier 162-5.}

Le tapuscrit conjecture ensuite (« je pense que ce résultat (facile) doit
pouvoir se généraliser ») un énoncé pour les topos. Soit $f : X \to X'$ un
morphisme de topos qui induit des isomorphismes en cohomologie pour tout
faisceau localement constant sur $X'$, le cas non commutatif inclus. Supposons
$X$ et $X'$ \emph{localement homotopiquement triviaux} : pour tout entier $n
\geq 1$, tout objet $U$ a un recouvrement $(U_i \to U)$ tel que
\begin{itemize}
  \item[a)] tout système local et toute section sur $U$ deviennent constants
  sur les $U_i$ ;
  \item[b)] pour tout groupe abélien $G$, les $H^{j}(U, G) \to H^{j}(U_i, G)$
  sont nuls pour $1 \leq j \leq n$.
\end{itemize}
Alors $f^{*} : D^{b}_{\mathrm{lc}}(X') \to D^{b}_{\mathrm{lc}}(X)$ serait une
équivalence. Il en irait de même avec un système local d'anneaux sur $X'$, et
$f$ induirait une équivalence entre coefficients locaux sur $X$ et sur
$X'$.\footnote{« $X$ et $X'$ » est écrit à la main au-dessus de « $X$ (resp.\
$X'$) » biffé, de sorte que « soit » et « trivial » restent au singulier ;
« $\geq 1$ » est ajouté à la main et « de $X'$ » biffé après « tout objet
$U$ ». La flèche de « $f : X \to X'$ » manque à la frappe. En tête de la page
4, la phrase a été reprise à la machine en plusieurs couches : « pleinement
fidèle (resp.\ » y est biffé, et « foncteur pl.\ fid.\ (resp.\ une
équivalence) » surcharge « isomorphisme » ; « une équivalence » est la lecture
retenue. Une parenthèse biffée à la machine se demandait si cette propriété ne
devait pas être mise dans les hypothèses sur $f$. Dans la phrase sur les
coefficients locaux, la page écrit $C$, $C'$ là où le contexte demande $X$,
$X'$.} Une note manuscrite au pied de la page 3, appelée en regard de la
définition, en précise la portée : la condition « n'est typiquement \emph{pas}
satisfaite par les schémas avec leur topologie étale (mais bien par […] loc.\
noeth.\ avec top.\ Zariski) ».\footnote{Le mot après « bien par » est
illisible sur 111, précédé d'un mot biffé ; la transcription de l'exemplaire
annoté (dossier 162-5, page 34) y lit, sous doute, « variétés ».
« typiquement », « avec » et « étale » sont aussi des lectures incertaines.
Le constat se vérifie au moins sur un exemple, et l'argument est le nôtre :
pour une courbe complexe $X$ de genre $\geq 1$ et un revêtement fini étale $U
\to X$ de degré $d$, la trace montre que $H^{1}(X, \mathbb{Z}/m) \to H^{1}(U,
\mathbb{Z}/m)$ est injectif pour $m$ premier à $d$. Aucun tel revêtement ne tue
donc à la fois tous les $H^{1}$, comme le demande b). Nous ne développons pas
le cas d'un recouvrement étale quelconque.}

Le principe de démonstration proposé est le suivant. Un topos localement
homotopiquement trivial est localement connexe et localement simplement
connexe, d'où une bonne théorie des $\pi_0$ et $\pi_1$, qui sont ici discrets.
On est ramené à un cas particulier du critère d'Artin et Mazur : un
homomorphisme de groupes $G \to H$ est un isomorphisme s'il induit des
bijections sur $H^{0}$ et $H^{1}$, cas non commutatif compris. La pleine
fidélité se ramène, par la suite spectrale, aux $\mathrm{Ext}^{i}$ de
coefficients locaux. Elle résulterait de deux faits : dans un topos localement
homotopiquement trivial, les faisceaux abéliens localement constants sont
stables par $\mathrm{Ext}^{i}$, et $M \mapsto M_X$ commute à ces
$\mathrm{Ext}^{i}$. Pour $X$, $X'$ localement homotopiquement triviaux, trois
conditions sur $f$ seraient alors équivalentes.
\begin{itemize}
  \item[a)] $f$ est un homotopisme : il induit un isomorphisme sur les $H^{i}$
  pour tout système local sur $X'$, non nécessairement commutatif.
  \item[b)] $f^{*} : D^{b}_{\mathrm{lc}}(X') \to D^{b}_{\mathrm{lc}}(X)$ est une
  équivalence.
  \item[a')] $f$ induit une équivalence entre systèmes locaux abéliens sur $X'$
  et sur $X$, et des isomorphismes sur les $H^{i}$ correspondants.
\end{itemize}
La page ignore si l'on peut, dans a), se borner aux systèmes locaux
commutatifs, et la main met en marge un point d'interrogation. La réponse est
non.\footnote{Remarque nôtre. Il existe des groupes acycliques non triviaux,
par exemple le groupe de Higman à quatre générateurs. Pour un tel groupe $G$,
le morphisme du topos $BG$ vers le point, entre topos localement
homotopiquement triviaux, induit des isomorphismes sur la cohomologie de tout
système local abélien du point. Ce n'est pas un homotopisme : $H^{1}(-, G)$ le
distingue, et a') échoue. Le phénomène est celui qu'exploite la construction
« plus » de Quillen, postérieure à ces notes.}

L'équivalence de a) et b) motive une définition des types d'homotopie par la
catégorie $D^{b}_{\mathrm{lc}}(X)$.\footnote{L'exposant et l'indice de ce
dernier $D$ sont ajoutés à la main.} On la munirait de la sous-catégorie pleine
des systèmes locaux, du foncteur cohomologique à valeurs dans celle-ci et du
produit tensoriel, ce qui fait sortir de $D^{b}$ vers $D^{-}$ : « redactor
demerdetur ».\footnote{L'idée de décrire un type d'homotopie par des
coefficients reviendra chez lui, sous une tout autre forme, dans « À la
poursuite des champs » (1983), auprès duquel le dossier est classé. Nous
n'allons pas plus loin.}

\pagerange{5}{10}
\section*{II. $n$-catégories, catégories $n$-uples, Gr-catégories (pages 5 à 10)}

\pagerange{5}{6}
\subsection*{Catégories doubles et 2-catégories (pages 5 et 6)}

Le tapuscrit appelle « bi-catégorie » une catégorie dans $(\mathrm{Cat})$.
C'est la donnée d'une catégorie $C_0$, de foncteurs source et but $p_1, p_2 :
C_1 \rightrightarrows C_0$, d'un foncteur unité $C_0 \to C_1$ et d'une
composition $\pi : C_1 \times_{C_0} C_1 \to C_1$, avec les axiomes habituels.
On dit aujourd'hui \emph{catégorie double}.\footnote{Ce n'est pas une
bicatégorie au sens de Bénabou (1967), c'est-à-dire une 2-catégorie faible :
c'est la catégorie double d'Ehresmann (1963). Nous employons le nom moderne
pour éviter le contresens. La page dit « un triple » pour une donnée à cinq
termes et ne mentionne pas le foncteur unité, que nous ajoutons. Sur cette
page et la suivante, la main barre les exposants tapés $0$, $1$, $00$, …
et les récrit en indices ; nous suivons la correction.} On identifie
$(\mathrm{Cat})$ à la sous-catégorie pleine de $(\mathrm{sSet})$ formée des
ensembles simpliciaux qui satisfont la condition d'exactitude « familière » :
transformer sommes amalgamées en produits fibrés, soit $X_n \simeq X_1
\times_{X_0} \cdots \times_{X_0} X_1$, ce qu'on appelle aujourd'hui la
condition de Segal.\footnote{La page laisse la phrase en suspens, « des objets
qui … », et énonce la condition une ligne plus bas, pour les objets
bisimpliciaux ; nous la complétons par celle-ci. Le nom de Segal est
postérieur.} De même, les catégories doubles s'identifient aux ensembles
bisimpliciaux $(X_{ij})$ qui satisfont cette condition ligne par ligne et
colonne par colonne. Les catégories triples s'identifient à des objets
trisimpliciaux, et ainsi de suite.

\emph{Exemple.} Pour une catégorie $C$, posons $C_0 = C$ et $C_1 =
\mathrm{Fl}(C)$, la catégorie des flèches, avec les structures évidentes. Les
composantes simpliciales supérieures sont les catégories $\mathrm{Fl}^{n}(C)$
des suites de $n$ flèches, et l'ensemble bisimplicial correspondant a pour
composantes $\mathrm{Hom}([n] \times [m], C)$. Les objets de $X_{11}$ sont les
carrés commutatifs de $C$, qu'on peut lire de deux manières comme morphismes de
flèches. L'échange des lignes et des colonnes est une symétrie de la catégorie
des catégories doubles ; de même, le groupe symétrique $\mathfrak{S}_n$ opère
sur la catégorie des catégories $n$-uples.

Une \emph{2-catégorie} (stricte) est une catégorie double dont la catégorie
d'objets $C_0$ est discrète : en termes bisimpliciaux, un ensemble bisimplicial
dont la première ligne est un ensemble simplicial constant. Il paraît qu'on a
une interprétation analogue des $n$-catégories en termes d'ensembles
$n$-simpliciaux. Le tapuscrit note de demander à Quillen le rapport entre les
$m$-catégories dans $(n\text{-Cat})$ et les ensembles
$(m+n)$-simpliciaux.\footnote{Un double trait vertical manuscrit, en marge,
signale ces deux phrases. Suit, biffé à la machine : « En particulier entre
cette dernière notion et celle de ». C'est la voie que prendront, pour les
$n$-catégories faibles, Tamsamani et Simpson à la fin des années 1990.}

\pagerange{7}{7}
\subsection*{Groupes dans $(\mathrm{Cat})$ et catégories de Picard (page 7)}

Les groupes dans $(\mathrm{Cat})$ sont les groupes simpliciaux qui satisfont
la condition de Segal, et qui sont donc déterminés par $C_0$ et $C_1$. Par la
correspondance de Dold--Kan (le « tapis de Dold »), les groupes
\emph{abéliens} dans $(\mathrm{Cat})$ forment une catégorie équivalente à celle
des complexes de chaînes de longueur 1 :
\[
0 \longleftarrow C_0 \xleftarrow{\ d\ } C_1 \longleftarrow 0 .
\]
Le groupe abélien dans $(\mathrm{Cat})$ associé à un tel complexe est la
catégorie de Picard du complexe, « “ma” construction chérie » : ses objets sont
les éléments de $C_0$, et ses flèches de $x$ vers $y$ les $z \in C_1$ tels que
$dz = y - x$. On obtient ainsi, à isomorphisme près, tous les groupes abéliens
de $(\mathrm{Cat})$.\footnote{La notion de catégorie de Picard est la sienne.
Deligne en donnera l'exposé de référence en 1973 (SGA 4, exposé XVIII, §
1.4), avec l'équivalence entre catégories de Picard strictement commutatives
et complexes de longueur 1 à quasi-isomorphisme près.}

Complément important : toute Gr-catégorie, c'est-à-dire toute catégorie
monoïdale où objets et flèches sont inversibles, est équivalente à un vrai
groupe dans $(\mathrm{Cat})$. Le tapuscrit avait étendu l'énoncé au cas
commutatif, par deux parenthèses tapées en interligne ; la main les biffe, et
avec elles le mot « commutative ».\footnote{La biffure est fondée. Une
catégorie de Picard symétrique dont la symétrie $c_{x,x}$ n'est pas l'identité
ne peut pas être équivalente à un groupe abélien strict dans
$(\mathrm{Cat})$ : seules les catégories de Picard strictement commutatives le
sont. En marge, en regard de ces lignes, une note manuscrite écrite de biais
est barrée de traits obliques et ne se lit pas sur 111. Sur l'exemplaire
annoté (dossier 162-5, page 38), la marge porte en outre une note conservée
selon laquelle une « cat.\ de Picard commut. » ne se réalise que par un
$(\mathrm{Cat})$-groupe non commutatif ; la transcription de 111 ne la
rapporte pas. Le terme de Gr-catégorie, qu'on lit ici en 1968, est celui que
reprendra la thèse de Hoàng Xuân Sính (1975), préparée sous sa direction ; on
dit aujourd'hui aussi 2-groupes.} Une ligne manuscrite ajoute : quand on
« simplifie » modulo l'équivalence de Grab-catégories, cela revient à
travailler dans la catégorie \emph{dérivée} des complexes, dont les seuls
invariants sont $H_0$ et $H_1$, autrement dit $\pi_0$ et $\pi_1$.\footnote{Les
mots « simplifie », « Grab », « seuls invariants », $H_0$, $H_1$ sont pâlis
sur la photocopie, et un mot y est illisible ; l'exemplaire annoté (dossier
162-5, page 38) les lit clairement. Pour les objets, l'énoncé est exact parce
que $\mathrm{Ext}^{2}_{\mathbb{Z}} = 0$ : un complexe de longueur 1 est
quasi-isomorphe à $H_0 \oplus H_1[1]$. Les morphismes voient encore
$\mathrm{Ext}^{1}(H_0, H_1)$.}

Une note manuscrite encadrée, en tête de la page 7, avertit : « Attention à la
catégorie $\mathrm{Hom}_{\mathrm{grab}}$, qui s'identifie à la catégorie
définie par un $R\mathrm{Hom}_{\mathbb{Z}}(C_{\bullet}, C'_{\bullet})$, ce
point devant être élucidé ! ». Les « Grab » sont les groupes abéliens, et il
s'agit des homomorphismes entre les catégories de Picard $P(C_\bullet)$,
$P(C'_\bullet)$ de deux complexes de longueur 1. L'énoncé correct demande une
troncature : les foncteurs additifs de $P(C_\bullet)$ dans $P(C'_\bullet)$
forment la catégorie de Picard associée à $\tau^{\leq 0}
R\mathrm{Hom}(C_\bullet, C'_\bullet)$, en indexation
cohomologique.\footnote{La troncature est notre précision : $R\mathrm{Hom}$
entre complexes de longueur 1 a de la cohomologie en trois degrés, $-1$, $0$
et $1$, et une catégorie de Picard n'en retient que deux. C'est sous cette
forme que l'énoncé figure chez Deligne (SGA 4 XVIII).}

On a ainsi une classification beaucoup plus restreinte que celle des H-espaces
(resp.\ des H-espaces homotopiquement commutatifs) à homotopie près. La raison
en est qu'on est plus exigeant, dans le cadre catégorique, sur l'associativité
(resp.\ la commutativité). On ne demande pas seulement l'existence
d'isomorphismes d'associativité, qui suffirait pour avoir un groupe dans
$\mathrm{Ho}(\mathrm{Cat}) \simeq \mathrm{Ho}(\mathrm{sSet})$ ; on les suppose
\emph{donnés}, avec des compatibilités, et c'est ce qui permet de se ramener à
de vrais groupes dans $(\mathrm{Cat})$. La main biffe la suite, « resp.\ de
vraie commutativité », pour la raison qu'on vient de dire.\footnote{« donnés »
est souligné à la main ; une croix manuscrite de renvoi suit
« d'associativité », et un crochet manuscrit entoure les deux dernières lignes.}

\pagerange{8}{9}
\subsection*{Le cran suivant : groupes dans $(2\text{-Cat})$ (pages 8 et 9)}

On doit trouver une meilleure approximation, d'un cran, de la classification
topologique, en passant aux groupes (resp.\ groupes abéliens) dans
$(2\text{-Cat})$, ou plus généralement aux « Gr-2-catégories » (resp.\
« Grab-2-catégories »). Deux opérations relient les crans. Une Grab-catégorie
définit une Grab-2-catégorie, comme une catégorie définit une 2-catégorie :
c'est l'analogue d'un « cosquelette tronqué ». Inversement, une
Grab-2-catégorie définit une Grab-catégorie, celle des automorphismes de
l'objet unité ; le tapuscrit y voit l'analogue du « tuage des groupes
d'homotopie en dimension $> 1$ ».\footnote{Remarque nôtre : à calculer les
invariants, la seconde opération ressemble plutôt au passage à l'espace des
lacets. Les $\pi_0$, $\pi_1$ de la Gr-catégorie des automorphismes de l'unité
sont les $\pi_1$, $\pi_2$ de la Gr-2-catégorie de départ ; c'est un décalage
plus qu'une troncature. Deux mots de la parenthèse portent un signe manuscrit
de renvoi.} Pour la notion « grossière », l'idée de sa lettre à Verdier était
qu'on doit trouver la même chose que les tronqués au cran 2 dans la catégorie
\emph{homotopique} des H-espaces, c'est-à-dire ce qu'on en déduit en tuant les
groupes d'homotopie en dimension $> 2$.\footnote{Les signes $>$ de ce passage
sont récrits à la main. La lettre à Verdier n'est pas dans le dossier.}

Le tapuscrit note de demander à Quillen si les objets-groupes et les objets
connexes de la catégorie homotopique pointée forment des catégories
équivalentes par le foncteur espace des lacets. La main répond en marge :
« non, on a un foncteur pl.\ fidèle (?) mais pas ess.\ surjectif ». La
non-surjectivité est juste, et le point d'interrogation justifié : le foncteur
des lacets n'est même pas fidèle.\footnote{Remarque nôtre. La classe $ab \in
H^{4}(S^{2} \times S^{2}; \mathbb{Z})$, produit des deux générateurs, définit
une application non homotope à zéro $S^{2} \times S^{2} \to K(\mathbb{Z}, 4)$ ;
sa bouclée est pourtant homotope à zéro, parce que la suspension
cohomologique annule les produits. La note de marge est écrite de biais ; la
question est signalée par un double trait vertical, et « ponctuée », tapé en
interligne, est mis à sa place par un trait.}

\emph{Les vrais groupes abéliens dans $(2\text{-Cat})$.} L'approche de Quillen
les identifie d'abord à des groupes abéliens dans la catégorie des catégories
doubles, donc à des groupes abéliens bisimpliciaux particuliers, soumis à la
condition de Segal sur chaque ligne et chaque colonne. Par la version
bisimpliciale de Dold--Kan, ils équivalent aux bicomplexes de groupes abéliens
nuls hors du carré $0 \leq i, j \leq 1$. La condition que la première ligne
corresponde à un groupe simplicial constant signifie qu'elle se réduit à son
terme de degré 0, et le bicomplexe devient\footnote{La frappe portait
« colonne », que la main biffe et remplace par « cond. ». Les deux
« $\leftarrow 0$ » de droite sont ajoutés à la main.}

\begin{tikzcd}
C_{00} & 0 \arrow[l] & 0 \arrow[l] \\
C_{10} \arrow[u] & C_{11} \arrow[l] \arrow[u] & 0 \arrow[l]
\end{tikzcd}

Ces bicomplexes s'identifient aux complexes de chaînes de longueur 2,
$C_{11} \to C_{10} \to C_{00}$, et la catégorie des groupes abéliens dans
$(2\text{-Cat})$ est donc équivalente à celle des complexes de longueur 2. La
2-catégorie associée à un tel complexe est « mon interprétation chérie ». Il
resterait à décrire, par un complexe $R\mathrm{Hom}$, la 2-Grab-catégorie des
homomorphismes entre deux telles 2-Grab-catégories strictes.

L'idée qui se dégage, et que Quillen, « on espère », va expliciter, est que
les groupes abéliens dans $(n\text{-Cat})$ correspondent aux complexes de
longueur $n$. On raisonnerait soit à isomorphisme près des deux côtés, soit à
$n$-équivalence près d'un côté et dans la catégorie dérivée de l'autre, au
moyen de $\mathrm{Hom}^{\bullet}$ (resp.\
$R\mathrm{Hom}^{\bullet}$).\footnote{Pour des groupes abéliens
\emph{stricts}, la catégorie dérivée est bien la bonne cible. Pour les
versions faibles et symétriques, elle ne l'est plus : ce sont des spectres
tronqués qui les classent, et non des complexes. La biffure « resp.\ de vraie
commutativité » de la page 7 enregistre déjà ce phénomène au cran 1.}

\pagerange{9}{10}
\subsection*{Les $K_n$ d'une catégorie monoïdale (pages 9 et 10)}

Soit $C$ une catégorie munie d'un produit tensoriel associatif, et
éventuellement unitaire : un bifoncteur avec des isomorphismes
d'associativité, mais sans inverses. On doit pouvoir lui associer une
Gr-catégorie universelle, comme on associe à un monoïde son groupe symétrisé,
d'où des invariants $K_0(C)$ et $K_1(C)$. On suppose qu'on trouvera de même,
pour tout $n$, une Gr-$n$-catégorie universelle où envoyer $C$ de façon
compatible à $\otimes$ ; ses invariants $\pi_n$ seraient par définition les
$K_n(C)$. Un cas intéressant est celui d'une catégorie à produits finis,
additive par exemple, où $\otimes$ est le produit. Mais pour une catégorie
abélienne il y a une notion plus adéquate du $K_0$, qui devrait correspondre à
une meilleure notion des $K_i$, « inspirée du tapis de Roos ». C'est sans doute
en termes de « vraies » catégories triangulées qu'il faudra poser les
définitions.\footnote{Un crochet ouvrant manuscrit précède « Mais on
notera ». Nous n'identifions pas le « tapis de Roos ». Quillen définira les
$K_n$ supérieurs en 1970--1972, par la construction « plus » puis par la
construction $Q$ (publiée en 1973) ; pour une catégorie monoïdale symétrique,
ce sera par complétion en groupe (Grayson, 1976, d'après Quillen).
Schlichting (2002) montrera que la $K$-théorie n'est pas un invariant de la
catégorie triangulée seule, ce qui donne du poids aux guillemets de
« vraies ».}

Peut-être Quillen a-t-il une définition simpliciale plus directe du H-espace
qui devrait correspondre à $C$, sans passer par les $n$-catégories. La marge
répond : « oui ».\footnote{Avec un double trait vertical. C'est en effet par un
espace que Quillen définira les $K_n$. La marge n'est pas datée.} Ce n'est
certainement pas en général le nerf $N(C)$ : une catégorie additive est
contractile, puisque le foncteur identique est homotope au foncteur nul, une
transformation naturelle de l'un à l'autre suffisant. Pourtant, lorsque $C$
est une Gr-catégorie, le nerf naïf est le bon : on trouve des $\pi_0$ et
$\pi_1$ qui sont aussi ceux de $C$ au sens des catégories sans Gr-structure.

Une note manuscrite sous la frappe ajoute : « A expliquer : pour des
Gr-$n$-catégories, les $\pi_i$ au sens semi-simplicial sont aussi les $\pi_i$
au sens Gr-$n$-catégories, savoir les objets à isom.\ près dans divers
$\mathrm{Hom}(1, 1)$ … ».\footnote{La seconde moitié de la note est barrée de
traits obliques, sous lesquels elle reste lisible par endroits. Elle parlait
d'une interprétation de l'opération des lacets, $\Omega$, en termes du
« shift géométrique » sur les $n$-catégories ; plusieurs mots en sont
incertains, et un est illisible. L'exemplaire annoté (dossier 162-5, page 41)
porte la même phrase.}

\pagerange{11}{16}
\section*{III. Point de vue « motivique » en théorie du cobordisme (pages 11 à 16)}

\pagerange{11}{13}
\subsection*{Le problème universel et la catégorie $\mathbf{B}$ (pages 11 à 13)}

\emph{Le problème.} Soit $\mathcal{V}$ la catégorie des variétés
différentiables compactes sans bord, non nécessairement orientables, dont les
morphismes sont les classes d'homotopie d'applications continues.\footnote{La
frappe ne dit pas « compactes ». Nous l'ajoutons : les images directes, la
composition des correspondances et la surjectivité de Thom invoquée plus bas
l'exigent. La main écrit justement « (compactes ?) » à cet endroit ; c'est
l'une des notes propres à la photocopie, absente de l'exemplaire annoté. La
généralisation de la page 16 lève ensuite l'hypothèse.} Soit $\mathcal{E}$ une
catégorie. On s'intéresse aux couples $(F_\bullet, F^\bullet)$ formés d'un
foncteur covariant $F_\bullet : \mathcal{V} \to \mathcal{E}$ et d'un foncteur
contravariant $F^\bullet : \mathcal{V}^{\circ} \to \mathcal{E}$, qui ont les
mêmes valeurs sur les objets, et qui satisfont le changement de base pour tout
carré cartésien d'applications différentiables où $f$ et $g$ sont transverses :

\begin{tikzcd}
X' \arrow[r, "g'"] \arrow[d, "f'"'] & X \arrow[d, "f"] \\
Y' \arrow[r, "g"'] & Y
\end{tikzcd}

\[
g^{\bullet} f_{\bullet} = f'_{\bullet}\, g'^{\bullet}
\qquad (f_\bullet = F_\bullet(f),\ f^\bullet = F^\bullet(f)).
\]
La question est de trouver un couple universel, et Quillen prouve que la
construction suivante en donne un.\footnote{Ce cadre, deux fonctorialités liées
par le changement de base transverse, est celui que Quillen rendra célèbre en
1971 pour le cobordisme complexe (« Elementary proofs of some results of
cobordism theory using Steenrod operations ») ; on le retrouve dans les
théories bivariantes de Fulton et MacPherson (1981) et dans le cobordisme
algébrique de Levine et Morel (2007). Tout cela est postérieur à ces notes. La
page note $B$ la catégorie but ; nous écrivons $\mathcal{E}$, pour garder
$\mathbf{B}$ à la solution.}

\emph{L'anneau $B(X)$.} Pour $X$ dans $\mathcal{V}$, soit $B(X)$ l'ensemble des
classes de variétés $Z \to X$ au-dessus de $X$ pour la relation de cobordisme :
$Z' \to X$ et $Z'' \to X$ sont cobordantes s'il existe une variété à bord $W$ et
$W \to X$ tels que $\partial W = Z' \sqcup Z''$, les deux applications étant
induites. La somme est la réunion disjointe ; c'est une loi de groupe où tout
élément est d'ordre 2. Le produit est le produit fibré, après qu'on a rendu les
deux applications transverses par le théorème de Thom. L'anneau est gradué par
la dimension relative : $B_d(X)$ est formé des $Z \to X$ avec $\dim Z - \dim X
= d$. Pour $f : X \to Y$, l'image directe $f_\bullet : B(X) \to B(Y)$ est la
composition avec $f$, et l'image inverse $f^\bullet : B(Y) \to B(X)$ le produit
fibré transverse. $B$ est ainsi un foncteur en anneaux gradués contravariant,
et un foncteur en groupes covariant, et la formule de changement de base
entraîne la formule de projection habituelle.\footnote{La graduation (« NB
L'anneau est de plus gradué par la dimension relative … », en interligne, et
« (gradué) ») et la parenthèse qui rattache la formule de projection au
changement de base sont des ajouts manuscrits ; « gradué » est une lecture
incertaine. En termes d'aujourd'hui, $B_d(X)$ est le groupe de cobordisme non
orienté $MO^{-d}(X)$ d'Atiyah (1961) ; par la dualité de Poincaré--Atiyah,
c'est aussi le bordisme $\mathfrak{N}_{n+d}(X)$ de Thom, si $X$ est de
dimension $n$. La page dit « anneau de bordisme ».}

\emph{La catégorie $\mathbf{B}$.} Les constructions usuelles de la théorie des
correspondances donnent une catégorie $\mathbf{B}$, dont les objets sont ceux
de $\mathcal{V}$ et dont les morphismes sont donnés par
\[
\mathrm{Hom}_{\mathbf{B}}(X, Y) = B(X \times Y),
\qquad
\beta \circ \alpha = p_{13\bullet}\big(p_{12}^{\bullet}\alpha \cdot
p_{23}^{\bullet}\beta\big),
\]
où les $p_{ij}$ sont les projections de $X \times Y \times Z$.\footnote{La page
renvoie aux « constructions bien connues » ; la formule de composition est la
nôtre, et c'est elle qui demande $Y$ compact.} La symétrie $B(X \times Y)
\simeq B(Y \times X)$ donne un isomorphisme canonique $\mathbf{B} \simeq
\mathbf{B}^{\circ}$, qui est l'identité sur les objets. Le foncteur $F_\bullet
: \mathcal{V} \to \mathbf{B}$ est l'identité sur les objets et envoie une
flèche sur son graphe ; composé avec la symétrie, il donne $F^\bullet :
\mathcal{V}^{\circ} \to \mathbf{B}$.\footnote{La frappe écrit « $F^{\bullet} :
C \to \mathbf{B}$ », sans l'opposée qu'appelle un foncteur contravariant.} Le
couple satisfait aux conditions voulues ; la page note que cela devrait être
formalisé pour des correspondances générales.

\emph{La propriété universelle.} Soit $(F'_\bullet, F'^\bullet)$ un couple à
valeurs dans une catégorie $\mathcal{E}'$, et notons $X' = F'(X)$. Il faut
définir des applications $B(X \times Y) \to \mathrm{Hom}_{\mathcal{E}'}(X',
Y')$. On envoie $Z \to X \times Y$, de composantes $p$ et $q$, sur le composé
\[
X' \xrightarrow{\ F'^{\bullet}(p)\ } Z' \xrightarrow{\ F'_{\bullet}(q)\ } Y' .
\]
Il faut voir que le résultat ne dépend que de la classe de cobordisme de $Z$
au-dessus de $X \times Y$.\footnote{La page écrit « un $Z_1$ cobordant au
dessus de $X$ » ; il s'agit d'un cobordisme au-dessus de $X \times Y$, comme
le dit la phrase suivante. « est le même » est ajouté à la main.} Soit $T \to
X \times Y$ un cobordisme, avec $\partial T = Z_0 \sqcup Z_1$. En « doublant »
$T$, on obtient une variété $W$ au-dessus de $S^{1}$, par $p : W \to S^{1}$,
transverse au-dessus de deux points $s_0$, $s_1$ dont les fibres sont $Z_0$ et
$Z_1$, et l'application vers $X \times Y$ se prolonge à $W$. Notant $\alpha_i :
Z_i \hookrightarrow W$ l'inclusion, on est ramené à montrer que le composé
\[
W' \xrightarrow{\ \alpha_i^{\bullet}\ } Z'_i \xrightarrow{\ \alpha_{i\bullet}\ } W'
\]
ne dépend pas de $i$. Considérons le carré cartésien d'inclusions, transverse
pour $s$ valeur régulière de $p$ :

\begin{tikzcd}
W_{s} \arrow[r, "\alpha_{s}"] \arrow[d, "\alpha_{s}"'] & W \arrow[d, "{\Gamma = (p, \mathrm{id}_{W})}"] \\
W \arrow[r, "i_{s}"'] & S^{1} \times W
\end{tikzcd}

Le changement de base donne
\[
\alpha_{s\bullet}\,\alpha_s^{\bullet} = \Gamma^{\bullet}\, i_{s\bullet},
\]
et le second membre ne dépend pas de $s$, puisque la classe d'homotopie de
$i_s$ n'en dépend pas. La page conclut : « Le reste est sans doute
l'AQT ».\footnote{La page écrit $X$ pour notre $W$, qui est son
$\overline{T}'$, et note $\overline{g}$ le prolongement. Elle écrit
$\alpha_s^{*}\alpha_{s*} = \Gamma^{*} i_{s*}$ : lu comme une composition, le
premier membre ne se compose pas, et nous écrivons l'ordre qui convient. Le
changement de base pris dans l'autre sens donne aussi $i_s^{\bullet}
\Gamma_{\bullet}$, qui convient de même. Les $\alpha$, les astérisques et les
étiquettes du carré sont portés à la main dans des blancs de la frappe ; la
première étiquette du composé se lit $\alpha_i$ sans astérisque, ici comme sur
l'exemplaire annoté, et l'étiquette $\Gamma = (p, \mathrm{id})$, pâle sur la
photocopie, a son signe $=$ souligné. Nous ne savons pas résoudre
l'abréviation « AQT ».}

\pagerange{13}{14}
\subsection*{Structure des $B_\bullet(X)$, enveloppe de Karoubi, l'objet $\Lambda(-1)$ (pages 13 et 14)}

\emph{L'anneau des coefficients.} Les $B_\bullet(X)$ sont des algèbres graduées
sur $\Lambda = B_\bullet(\mathrm{pt})$. $\Lambda$ est à degrés positifs et
réduit à $\mathbb{F}_2$ en degré 0 : c'est l'anneau de bordisme non orienté
$\mathfrak{N}_*$ de Thom. C'est une algèbre de polynômes $\mathbb{F}_2[x_i]$,
avec un générateur $x_i$ de degré $i$ pour chaque entier $i \geq 1$ qui n'est
pas de la forme $2^{\nu} - 1$.\footnote{La description des générateurs est une
insertion manuscrite, pâle et rognée au bord droit sur 111, lue avec
l'exemplaire annoté (dossier 162-5, page 45) ; « entier $\geq 1$ » et « $\neq$
de la forme $2^{\nu} - 1$ » restent des lectures incertaines. Elle est
exacte.}

On a un homomorphisme évident d'anneaux gradués $B_\bullet(X) \to
H^{\bullet}(X)$. Il envoie $Z \to X$ sur l'image de $1 \in H^{0}(Z)$ par
l'homomorphisme de Gysin, envoie $B_{-k}$ dans $H^{k}$, et il est donc nul en
degrés $> 0$. Il est surjectif par Thom : toute classe provient d'une variété
compacte.\footnote{$H^{\bullet}(X)$ est la cohomologie à coefficients dans
$\mathbb{F}_2$, comme le précise en marge une note propre à la photocopie :
« NB $H^{*}(X)$ $(= H^{*}(X, \mathbb{F}_2))$ ». La surjectivité combine le
théorème de représentabilité de Thom pour l'homologie modulo 2 et la dualité
de Poincaré, d'où, encore, l'hypothèse de compacité.} Une note manuscrite
ajoute : « en fait $H^{\bullet}(X) \simeq B_\bullet(X)/\Lambda^{+}B_\bullet(X)$
», et c'est exact.\footnote{« en fait » est une lecture incertaine.}

On prouve même qu'il existe un relèvement fonctoriel $s : H^{\bullet}(X) \to
B_\bullet(X)$ tel que l'homomorphisme induit
\[
H^{\bullet}(X) \otimes_{\mathbb{F}_2} \Lambda \xrightarrow{\ \sim\ }
B_\bullet(X)
\]
soit un isomorphisme de $\Lambda$-modules. Le tapuscrit ajoute que ce sera vrai
pour tout relèvement pris individuellement ; c'est exact pour tout relèvement
homogène.\footnote{Justification (nôtre). $B_\bullet(X)$ est un
$\Lambda$-module libre (Conner et Floyd, 1964), de quotient $H^{\bullet}(X)$
modulo $\Lambda^{+}$. Par le lemme de Nakayama gradué, les images d'une base
homogène engendrent, et un épimorphisme entre modules libres de même rang
gradué, de dimension finie en chaque degré, est un isomorphisme.} Le tapuscrit
disait le relèvement « compatible avec multiplication » ; la main met ces mots
entre parenthèses et écrit au-dessus « il est douteux qu'il soit ». Elle fait
aussi de la seconde condition une question : « Quillen ignore si » le
relèvement est compatible avec les homomorphismes de Gysin.\footnote{« a) »
est raturé à la main ; devant b), un crochet manuscrit ouvre la clause, et
« Quillen ignore si » est écrit au-dessus de « Le » biffé. Un point
d'interrogation et un grand trait en arc sont en marge. Ces deux questions,
multiplicativité et compatibilité aux images directes, sont exactement celles
que tranchera la note de Quillen de 1969, « On the formal group laws of
unoriented and complex cobordism theory ». La loi de groupe formel du
cobordisme non orienté y est universelle parmi les lois $F$ sur les
$\mathbb{F}_2$-algèbres telles que $F(x,x) = 0$ ; toute telle loi est
isomorphe à la loi additive sans lui être égale, et c'est cet écart qui
gouverne la compatibilité d'un relèvement multiplicatif aux images de Gysin.
Ces notes de septembre 1968 ne pouvaient pas le savoir, et nous ne leur
prêtons rien de ce qui précède.}

\emph{Modules et enveloppe de Karoubi.} Le foncteur $X \mapsto B_\bullet(X)$
est pleinement fidèle, de $\mathbf{B}$ dans la catégorie des $\Lambda$-modules
gradués libres de type fini. Les modules obtenus ont leurs générateurs en
degrés $\leq 0$ ; pour $X$ connexe, il y a un seul générateur de degré 0, et ils
satisfont la dualité de Poincaré.\footnote{La pleine fidélité résulte de
l'isomorphisme précédent, de la formule de Künneth et de la dualité de
Poincaré ; la page ne la détaille pas.} Passons à l'enveloppe pseudo-abélienne
$\widetilde{\mathbf{B}}$ de $\mathbf{B}$, c'est-à-dire à son enveloppe de
Karoubi, puisque $\mathbf{B}$ est déjà additive. On obtient tous les
$\Lambda$-modules gradués libres de type fini à générateurs en degrés $\leq 0$.
C'est une catégorie additive et karoubienne, munie d'un $\otimes$ induit par le
produit des variétés, qui a toutes les vertus, sauf d'être abélienne et d'avoir
des $\mathrm{Hom}$ internes.\footnote{La page dit « catégorie
quasi-abélienne », au sens de pseudo-abélienne. Ce n'est pas la catégorie
quasi-abélienne au sens moderne (Schneiders, 1999), notion toute différente.
Deux notes manuscrites propres à la photocopie le disent en toutes lettres :
« i.e.\ l'enveloppe “Karoubienne” », de biais en regard du mot, que la main
souligne d'un trait ondulé (page 13), et « $=$ karoubienne $+$ additive »,
en tête de la page 14. Le tilde de $\widetilde{\mathbf{B}}$ est ajouté à la
main.}

\emph{L'objet $\Lambda(-1)$.} Pour obtenir les $\mathrm{Hom}$ internes, on peut
paraphraser la construction des motifs. Soit $p : S^{1} \to e$ l'application
du cercle sur un point. Le morphisme $p^{\bullet} : e \to S^{1}$ de
$\mathbf{B}$ est scindé par le choix d'un point de $S^{1}$, et l'on pose
\[
\Lambda(-1) = \mathrm{coker}\big(p^{\bullet} : e \to S^{1}\big)
\ \in\ \widetilde{\mathbf{B}},
\qquad
B_\bullet(S^{1}) = \Lambda \oplus \Lambda(-1).
\]
C'est un module libre sur un générateur de degré $-1$, donc un objet de poids
1, le poids étant le degré des générateurs changé de signe.\footnote{La frappe
écrit $S$ pour le cercle, et « $\Lambda(-1)$ » est ajouté à la main, comme
« en » devant « considérant ». Le scindage, qui fait exister le conoyau dans
l'enveloppe karoubienne, est notre précision. C'est l'analogue de la
décomposition du motif de la droite projective en l'unité et le motif de
Lefschetz.} Le foncteur $M \mapsto M(-1) = M \otimes \Lambda(-1)$ de
$\widetilde{\mathbf{B}}$ dans elle-même est pleinement fidèle, et une
construction formelle standard permet de rendre $\Lambda(-1)$ inversible. Dans
la catégorie obtenue, les $\mathrm{Hom}$ internes existent et valent $\check{X}
\otimes Y$ ; on y trouve \emph{tous} les $\Lambda$-modules gradués libres de
type fini.\footnote{Le tilde et le signe de $\check{X}$ sont portés à la
main.}

Cette catégorie n'est pas abélienne, puisque $\Lambda$ n'est pas un corps. On
s'attend à ce que la catégorie abélienne engendrée soit celle des
$\Lambda$-modules gradués de présentation finie. Elle est bien abélienne,
quoique $\Lambda$ ne soit pas noethérien : $\Lambda$ est limite inductive
d'anneaux de polynômes en un nombre fini de variables, à morphismes de
transition plats, et il est donc cohérent.\footnote{Le tapuscrit demandait si
$\Lambda$ était noethérien, et si le nombre de générateurs était fini ; la
main biffe la question et répond par la phrase que nous rendons. « de type »
est corrigé en « de présentation » à la main. Le « lim » manuscrit ne porte pas
de flèche visible. Le terme d'anneau cohérent est le nôtre.}

\pagerange{14}{16}
\subsection*{L'analogie avec les motifs, la catégorie $\mathbf{S}$ et le groupe de Steenrod (pages 14 à 16)}

\emph{L'analogie.} Cette construction ressemble à une théorie des motifs où
l'on n'aurait pas passé à l'équivalence homologique des cycles, mais gardé
l'équivalence rationnelle, bien trop fine à beaucoup d'égards.\footnote{On
parle aujourd'hui de motifs de Chow, dans le cadre exposé par Manin en 1968.
« à beaucoup » et « convient », un peu plus loin, sont des lectures que la
transcription a dû établir sur des frappes surchargées ; la seconde reste
incertaine.} Pour un équivalent plus étroit, on pourrait poser le même
problème universel en prenant l'équivalence homologique des morphismes, via le
graphe, au lieu de l'équivalence homotopique. Un passage biffé à la main en
tirait une conclusion décevante : la catégorie solution aurait pour flèches
toutes les correspondances cohomologiques, et donnerait, après inversion de
$\mathbb{F}_2(-1)$, tous les espaces vectoriels gradués de dimension finie sur
$\mathbb{F}_2$, « un peu trop trivial comme structure ».\footnote{Le passage
commence à la dernière ligne de la page 14 ; il est biffé d'un trait par ligne
sur la page 14 et de grands traits en zigzag sur la page 15.}

\emph{La catégorie $\mathbf{S}$.} Pour trouver une catégorie universelle, on
associe à $X$ l'ensemble $S(X) \subset H^{\bullet}(X)$ des classes de
sous-variétés de $X$, modulo équivalence homologique. On prouve, par un
théorème de Thom dont il va être question, que c'est un sous-anneau ; par le
même théorème, les $S(X \times Y)$ sont stables par composition des
correspondances cohomologiques, ce qui ne semble pas trivial, puisque les
$S(X)$ ne sont pas stables par Gysin. D'où une catégorie additive
$\mathbf{S}$, dont les objets sont les variétés et les flèches les éléments des
$S(X \times Y)$. « Il semble qu'elle doive jouer le rôle universel que j'ai
envisagé plus haut. »\footnote{« él.\ des » est ajouté à la main.}

\emph{Le groupe $G$.} Pour en donner une interprétation algébrique,
considérons l'algèbre de Steenrod $\mathcal{A}$ des opérations cohomologiques
stables modulo 2. Elle est graduée, et munie d'une diagonale associative et
unitaire, anticommutative, c'est-à-dire commutative (on est en
caractéristique 2). Son dual gradué $\mathcal{A}_*$ est une algèbre
commutative, avec une diagonale associative et unitaire, et réduite à
$\mathbb{F}_2$ en degré 0. C'est donc l'anneau d'un schéma en groupes affine
$G$ sur $\mathbb{F}_2$. La page l'identifie au foncteur des
$\otimes$-automorphismes additifs de $H^{\bullet}$ sur $\mathcal{V}$ qui
opèrent trivialement sur $H^{1}(S^{1})$.\footnote{Après « stables », une
définition tapée de l'anneau des opérations est biffée à la main ; un long
trait vertical et un signe de renvoi sont en marge. La restriction « qui
opèrent trivialement sur $H^{1}(S^{1})$ » est manuscrite ; elle biffe un mot
avant $H^{1}(S^{1})$ et se termine sur un mot raturé. C'est elle qui écarte les
automorphismes de graduation. Aujourd'hui, $\mathcal{A}_* = \mathbb{F}_2[\xi_1,
\xi_2, \ldots]$ (Milnor, 1958), et $G(R)$ est le groupe, pour la composition,
des séries $x + \sum a_i x^{2^{i}}$ : les automorphismes stricts du groupe
formel additif.}

Le tapuscrit énonce alors, comme « le théorème de Thom », qu'un $x \in
H^{\bullet}(X)$ est dans $S(X)$ si et seulement s'il est fixe par $G$,
c'est-à-dire si $\mathcal{A}^{+} x = 0$, et demande aussitôt : « Pourquoi
cette condition est-elle déjà nécessaire ? ». Nous rapportons cet énoncé sous
réserve : lu avec l'action usuelle des carrés de Steenrod sur la cohomologie,
il est en défaut.\footnote{Remarque nôtre. Une droite $\mathbb{RP}^{1} \subset
\mathbb{RP}^{2}$ est une sous-variété, de classe le générateur $a \in
H^{1}(\mathbb{RP}^{2})$, et $\mathrm{Sq}^{1}a = a^{2} \neq 0$. La nécessité
échoue donc, comme la page le soupçonne. L'énoncé doit viser une autre action,
par exemple sur l'homologie, ou tordue par les classes de Wu, que la page ne
permet pas de restituer. Ce qui en dépend ci-dessous est rapporté sous la
même réserve.} Sous cette réserve, le foncteur $H^{\bullet}$, de $\mathbf{S}$
vers les représentations de $G$ dans des espaces vectoriels sur $\mathbb{F}_2$,
serait pleinement fidèle. Son enveloppe pseudo-abélienne serait une
sous-catégorie pleine des représentations de $G$ ; il est douteux qu'elle soit
abélienne, et plus encore qu'elle contienne toutes les représentations de
dimension finie.

$G$ est prolisse et, « d'après Quillen », pro-unipotent. Une note manuscrite
en marge confirme : « oui, car son anneau affine est une algèbre de polynômes
à une infinité d'indéterminées », réunion de sous-algèbres de type fini
stables par la diagonale $\Delta$.\footnote{Note de sept lignes écrite de biais
dans la marge gauche, rattachée par une flèche à la parenthèse sur $G$ ; le
premier mot est rogné, plusieurs mots sont incertains et trois sont illisibles
sur 111. L'exemplaire annoté (dossier 162-5, page 48) la porte plus lisible ;
nous en donnons le sens. Le « car » justifie la prolissité, par les
sous-algèbres $\mathbb{F}_2[\xi_1, \ldots, \xi_n]$ stables par la diagonale ;
la pro-unipotence vient de ce que l'algèbre de Hopf est graduée connexe,
réduite à $\mathbb{F}_2$ en degré 0.} Ici $G$ joue le rôle d'un \emph{groupe
de Galois motivique}. Contrairement à la situation familière, il est fortement
non réductif, et son corps de définition est $\mathbb{F}_2$, non
$\mathbb{Q}$.\footnote{L'expression « groupe de Galois motivique » est sur la
page, en 1968. Le formalisme tannakien qui lui donne son sens précis sera
publié par Saavedra Rivano en 1972. Nous ne tirons de cette rencontre aucune
conclusion de priorité.}

\pagerange{16}{16}
\subsection*{Orientations (page 16)}

Quillen généralise la construction de $\mathbf{B}$ de deux façons. D'abord aux
variétés non nécessairement compactes, $f_\bullet$ n'étant alors défini que
pour $f$ propre. Ensuite à une condition d'orientation stable sur les
\emph{morphismes}, donnée par un morphisme $H \to BO(\infty)$, où $H$ est un
H-espace dont la page se demande s'il est commutatif, associatif et unitaire.
Une orientation de $f : X \to Y$ est un relèvement homotopique à $H$ de
l'application $X \to BO$ qui classe $T_X - f^{\bullet}T_Y$. Les morphismes
$\mathrm{Hom}_{\mathbf{B}^{H}}(X, Y)$ sont les classes de diagrammes $X
\leftarrow Z \to Y$ de morphismes orientés, le second propre, modulo les bords
de tels diagrammes. Comme ensemble, cela se décrit par la variété $X \times Y$
munie de la famille des fermés propres sur $Y$.\footnote{La page note $H(X,Y)$
ces ensembles de morphismes, avec la même lettre que l'H-espace ; nous
changeons la notation. « stable » est tapé en interligne et rattaché à sa place
par un trait manuscrit ; le « t » de « Quillent » est barré à la main. Pour $H
= BO$ on retrouve $\mathbf{B}$ ; pour $H = BU$, on obtient le cobordisme
complexe $MU$, celui où Quillen trouvera en 1969 la loi de groupe formel
universelle. Ce sont les $(B,f)$-structures de Lashof (1963).}

Une note manuscrite au pied de la page ajoute : « Les $B^{H}_{m}(X)$ ($X$
compact disons) peuvent se décrire homotopiquement à la Thom, comme les
groupes d'homotopie stables d'espaces $\mathrm{Hom}(X, MH(N-m))$ … ». C'est la
construction de Pontryagin--Thom, $MH$ désignant l'espace de Thom associé à
$H$.\footnote{L'indexation est celle de la page, qui ne l'explique pas. La page
écrit un $\mathrm{Hom}$ souligné, c'est-à-dire un espace d'applications.}

\pagerange{1}{16}
\section*{La couche manuscrite de la photocopie, et sa date}

La couche manuscrite de 111 a deux strates, que la photocopie a fondues en un
seul gris. La première est l'encre bleue de l'exemplaire annoté du dossier
162-5 : la date, les corrections de notation, les réponses « oui » et « non,
… », les insertions des pages 11 à 16. La seconde est faite des marques que
porte aussi la photocopie de 162-5 et non l'exemplaire annoté. La
transcription de 111 en relève six : « (compactes ?) » (page 11), « NB
$H^{*}(X)$ » et « i.e.\ l'enveloppe “Karoubienne” » (page 13), « $=$
karoubienne $+$ additive » et le $\Lambda$ porté dans un blanc de la frappe
(page 14), et « cf exposé Karoubi à Bourbaki … » (page 16).\footnote{La lecture
modernisée du dossier 162-5 range aussi parmi les marques propres à sa
photocopie le « Si ! » de la page 3, « cond. » sur « colonne » (page 8), « en »
(page 14), « él.\ des » (page 15) et le t biffé de « Quillent » (page 16). La
transcription de 111 les porte sans les attribuer à l'une ou l'autre strate ;
nous ne tranchons pas. Le « au crayon » d'une première lecture de 111 n'est
pas confirmé : ces notes sont grises comme le reste.} Sur 111 comme sur la
photocopie de 162-5, ces marques ont les mêmes traits aux mêmes places : ce
sont des reproductions d'une seule écriture, non une écriture reprise.

Les deux premières précisent des hypothèses que le texte laisse implicites,
les deux suivantes traduisent son vocabulaire dans celui qui s'imposera, et le
$\Lambda$ comble un blanc. La
dernière, d'une encre plus pâle, pointe selon toute vraisemblance vers l'exposé
de Karoubi au Séminaire Bourbaki de 1971-1972, « Cobordisme et groupes formels
(d'après D. Quillen et T. tom Dieck) ». Cet exposé rend compte des résultats de
Quillen sur le cobordisme et les lois de groupes formels, ceux-là mêmes qui
répondent aux questions de la page 13. Si c'est bien lui, la note est
postérieure à 1971, et 111, qui la reproduit, a été tiré plus tard encore.
C'est une inférence, non une date lue. La datation du dossier reste celle de
l'inventaire, « 1968 », qui est celle du texte dactylographié ; et le
classement du dossier parmi les pièces rassemblées autour de « À la poursuite
des champs » est celui des archivistes, qui ne date pas la copie.

\end{document}
